\noindent\textbf{Supplementary Table S13A.} Halfvarson UC extent follow-up. Cases were extensive UC samples (E3) and controls were distal UC samples (E1/E2). The sample-level follow-up included 120 extensive and 173 distal samples from 25 and 35 subjects, respectively. Each row shows the top q-ranked dCST at the specified depth in rebuilt-cohort or direct AGP-derived transfer mode.

\begin{center}
\scriptsize
\begin{tabularx}{\textwidth}{@{} l r r Y Y r r Y @{}}
\toprule
Mode & Depth & Positive/negative samples & Top dCST & Frequentist OR (95\% CI) & \(p\) & \(q\) & Bayesian OR (95\% CrI) \\
\midrule
Rebuilt & 1 & 120/173 & \emph{Bacteroides dorei} & OR 0.554 [0.276, 1.1] & 0.072 & 0.072 & OR 0.553 [0.29, 1.05] \\
Rebuilt & 2 & 120/173 & \emph{Segatella sp.}\dcstsep \emph{Faecalibacterium prausnitzii} & OR 1.81 [0.908, 3.62] & 0.072 & 0.237 & OR 1.81 [0.963, 3.47] \\
Rebuilt & 3 & 120/173 & \emph{Segatella sp.}\dcstsep \emph{Faecalibacterium prausnitzii}\dcstsep \emph{Escherichia-Shigella coli} & OR 0 [0, 0.622] & 0.006 & 0.089 & OR 0.0304 [0.000063, 0.374] \\
Rebuilt & 4 & 120/173 & \emph{Segatella sp.}\dcstsep \emph{Faecalibacterium prausnitzii}\dcstsep \emph{Escherichia-Shigella coli} & OR 0 [0, 0.622] & 0.006 & 0.108 & OR 0.0303 [0.0000718, 0.38] \\
AGP transfer & 1 & 120/173 & \emph{Faecalibacterium prausnitzii} & OR 2.75 [1.26, 6.24] & 0.006 & 0.186 & OR 2.75 [1.33, 5.76] \\
AGP transfer & 2 & 120/173 & \emph{Bacteroides dorei}\dcstsep \emph{Akkermansia muciniphila} & OR inf [1.74, inf] & 0.004 & 0.333 & OR 41.7 [3, 19020] \\
AGP transfer & 3 & 120/173 & \emph{Bacteroides dorei}\dcstsep \emph{Akkermansia muciniphila}\dcstsep \emph{Faecalibacterium prausnitzii} & OR inf [1.74, inf] & 0.004 & 0.442 & OR 40.7 [3.02, 21142] \\
AGP transfer & 4 & 115/170 & \emph{Bacteroides dorei}\dcstsep \emph{Akkermansia muciniphila}\dcstsep \emph{Faecalibacterium prausnitzii} & OR inf [1.79, inf] & 0.004 & 0.407 & OR 41.9 [3, 20321] \\
\bottomrule
\end{tabularx}
\end{center}
\bigskip

\noindent\textbf{Supplementary Table S13B.} Halfvarson Crohn disease location follow-up. Cases were ileal-involved Crohn samples (Montreal L1/L1+L4 and L3/L3+L4) and controls were colonic-only Crohn samples (Montreal L2). The sample-level follow-up included 133 ileal-involved and 77 colonic-only samples from 31 and 18 subjects, respectively. Each row shows the top q-ranked dCST at the specified depth in rebuilt-cohort or direct AGP-derived transfer mode.

\begin{center}
\scriptsize
\begin{tabularx}{\textwidth}{@{} l r r Y Y r r Y @{}}
\toprule
Mode & Depth & Positive/negative samples & Top dCST & Frequentist OR (95\% CI) & \(p\) & \(q\) & Bayesian OR (95\% CrI) \\
\midrule
Rebuilt & 1 & 133/77 & \emph{Bacteroides dorei} & OR 1.59 [0.795, 3.18] & 0.182 & 0.182 & OR 1.6 [0.831, 3.04] \\
Rebuilt & 2 & 133/77 & \emph{Bacteroides dorei}\dcstsep \emph{Faecalibacterium prausnitzii} & OR 0.179 [0.0737, 0.407] & 7.3e-06 & 4.4e-05 & OR 0.178 [0.0805, 0.374] \\
Rebuilt & 3 & 133/77 & \emph{Bacteroides dorei}\dcstsep \emph{Faecalibacterium prausnitzii}\dcstsep \emph{Bacteroides stercoris} & OR 0.179 [0.0737, 0.407] & 7.3e-06 & 6.6e-05 & OR 0.179 [0.0796, 0.378] \\
Rebuilt & 4 & 133/77 & \emph{Bacteroides dorei}\dcstsep \emph{Faecalibacterium prausnitzii}\dcstsep \emph{Bacteroides stercoris}\dcstsep \emph{Agathobacter faecis} & OR 0.194 [0.0518, 0.611] & 0.002 & 0.019 & OR 0.198 [0.0612, 0.551] \\
AGP transfer & 1 & 133/77 & \emph{Faecalibacterium prausnitzii} & OR 0.119 [0.0279, 0.391] & 4.2e-05 & 0.001 & OR 0.121 [0.0355, 0.343] \\
AGP transfer & 2 & 129/77 & \emph{Bacteroides ovatus}\dcstsep \emph{Bacteroides dorei} & OR inf [1.58, inf] & 0.008 & 0.297 & OR 32.5 [2.59, 15243] \\
AGP transfer & 3 & 124/76 & \emph{Bifidobacterium breve}\dcstsep \emph{Bacteroides dorei} & OR 8.72 [1.26, 378] & 0.019 & 0.569 & OR 7.49 [1.61, 87.2] \\
AGP transfer & 4 & 122/75 & \emph{Bifidobacterium breve}\dcstsep \emph{Bacteroides dorei} & OR 8.76 [1.26, 379] & 0.019 & 0.574 & OR 7.73 [1.63, 94.5] \\
\bottomrule
\end{tabularx}
\end{center}
\bigskip

\noindent\textbf{Supplementary Table S13C.} Halfvarson Crohn disease calprotectin follow-up. Cases were Crohn samples with fecal calprotectin $\geq 250$ \textmu g/g and controls were Crohn samples with fecal calprotectin $< 250$ \textmu g/g; this was used as a conservative high-inflammatory-burden split rather than a remission cut-off. The sample-level follow-up included 103 high- and 99 low-calprotectin samples across 49 Crohn subjects; 13 subjects contributed both high- and low-calprotectin samples over time. Each row shows the top q-ranked dCST at the specified depth in rebuilt-cohort or direct AGP-derived transfer mode.

\begin{center}
\scriptsize
\begin{tabularx}{\textwidth}{@{} l r r Y Y r r Y @{}}
\toprule
Mode & Depth & Positive/negative samples & Top dCST & Frequentist OR (95\% CI) & \(p\) & \(q\) & Bayesian OR (95\% CrI) \\
\midrule
Rebuilt & 1 & 103/99 & \emph{Bacteroides dorei} & OR 1.71 [0.853, 3.5] & 0.139 & 0.139 & OR 1.72 [0.896, 3.33] \\
Rebuilt & 2 & 103/99 & \emph{Bacteroides dorei}\dcstsep \emph{Bacteroides uniformis} & OR 2.55 [1.05, 6.7] & 0.029 & 0.176 & OR 2.55 [1.15, 6.15] \\
Rebuilt & 3 & 103/99 & \emph{Segatella sp.}\dcstsep \emph{Faecalibacterium prausnitzii}\dcstsep \emph{Escherichia-Shigella coli} & OR 0.145 [0.0153, 0.677] & 0.005 & 0.044 & OR 0.154 [0.0284, 0.557] \\
Rebuilt & 4 & 103/99 & \emph{Segatella sp.}\dcstsep \emph{Faecalibacterium prausnitzii}\dcstsep \emph{Escherichia-Shigella coli} & OR 0.145 [0.0153, 0.677] & 0.005 & 0.054 & OR 0.153 [0.0257, 0.569] \\
AGP transfer & 1 & 103/99 & \emph{Bacteroides fragilis} & OR 0.242 [0.0419, 0.954] & 0.027 & 0.477 & OR 0.25 [0.0583, 0.791] \\
AGP transfer & 2 & 102/96 & \emph{Segatella sp.}\dcstsep \emph{Bacteroides dorei} & OR 0.0966 [0.00217, 0.721] & 0.008 & 0.456 & OR 0.111 [0.00896, 0.547] \\
AGP transfer & 3 & 100/94 & \emph{Segatella sp.}\dcstsep \emph{Bacteroides dorei}\dcstsep \emph{Faecalibacterium prausnitzii} & OR 0.11 [0.00243, 0.845] & 0.016 & 1.000 & OR 0.126 [0.0107, 0.635] \\
AGP transfer & 4 & 97/94 & \emph{Segatella sp.}\dcstsep \emph{Bacteroides dorei}\dcstsep \emph{Faecalibacterium prausnitzii}\dcstsep \emph{Agathobacter faecis} & OR 0 [0, 0.799] & 0.013 & 0.941 & OR 0.0338 [0.0000714, 0.466] \\
\bottomrule
\end{tabularx}
\end{center}
\bigskip

